%═══════════════════════════════════════════════════════════════════════════════
\chapter{Álgebra lineal numérica y espacios de funciones}
\label{ap:algebra_lineal}
%═══════════════════════════════════════════════════════════════════════════════

Este apéndice cubre las herramientas de álgebra lineal y análisis funcional que
el libro usa de forma recurrente pero que no siempre se tratan con suficiente
detalle en los cursos de física de pregrado. El contenido es estrictamente el
que aparece en el texto: SVD y pseudoinversa (Capítulos~\ref{ch:cap2},
\ref{ch:cap8}), normas matriciales (Capítulos~\ref{ch:cap3}, \ref{ch:cap8},
\ref{ch:cap9}), el hessiano y la curvatura del paisaje
(Capítulos~\ref{ch:cap3}, \ref{ch:cap5}), el espacio $L^2$ como espacio de
Hilbert (Capítulo~\ref{ch:cap10}), y los operadores adjuntos
(Capítulos~\ref{ch:cap6}, \ref{ch:cap11}).

%───────────────────────────────────────────────────────────────────────────────
\section{La descomposición en valores singulares (SVD)}
\label{sec:ap_svd}
%───────────────────────────────────────────────────────────────────────────────

\textit{Usado en: Capítulo~\ref{ch:cap2} (regresión regularizada, pseudoinversa,
curva L), Capítulo~\ref{ch:cap8} (doble descenso, filtro espectral de la parada
temprana).}

\medskip
La SVD es la factorización fundamental de cualquier matriz rectangular. A
diferencia de la diagonalización, existe para toda matriz sin restricción alguna.

\begin{teorema}{Descomposición en valores singulares}
\label{thm:ap_svd}
Sea $\bm{A}\in\mathbb{R}^{m\times n}$ con $\mathrm{rank}(\bm{A}) = r$. Existen
matrices ortogonales $\bm{U}\in\mathbb{R}^{m\times m}$, $\bm{V}\in\mathbb{R}^{n\times n}$
y una matriz diagonal $\bm{\Sigma}\in\mathbb{R}^{m\times n}$ tal que:
\begin{equation}
  \bm{A} = \bm{U}\bm{\Sigma}\bm{V}^\top,
  \qquad
  \bm{\Sigma} = \mathrm{diag}(\sigma_1,\ldots,\sigma_r,0,\ldots,0),
  \label{eq:ap_svd}
\end{equation}
con $\sigma_1 \geq \sigma_2 \geq \cdots \geq \sigma_r > 0$. Los escalares
$\{\sigma_i\}$ son los \textbf{valores singulares}; las columnas de $\bm{U}$
son los \textbf{vectores singulares izquierdos} $\{\bm{u}_i\}$ y las de
$\bm{V}$ los \textbf{vectores singulares derechos} $\{\bm{v}_i\}$.
\end{teorema}

La interpretación geométrica es la más útil para el libro: $\bm{A}$ mapea la
bola unitaria de $\mathbb{R}^n$ a un elipsoide en $\mathbb{R}^m$. Los vectores
singulares derechos $\{\bm{v}_i\}$ son los ejes del dominio; los izquierdos
$\{\bm{u}_i\}$ son los ejes del elipsoide imagen; y los valores singulares
$\{\sigma_i\}$ son las longitudes de los semiejes.

\medskip\noindent\textbf{Forma compacta.} La SVD se escribe en forma de suma de
rango 1:
\begin{equation}
  \bm{A} = \sum_{i=1}^r \sigma_i\,\bm{u}_i\bm{v}_i^\top.
  \label{eq:ap_svd_suma}
\end{equation}
Truncar la suma a los $k < r$ términos de mayor $\sigma_i$ produce la mejor
aproximación de rango $k$ de $\bm{A}$ en norma de Frobenius (teorema de
Eckart-Young).

\medskip\noindent\textbf{Cuatro subespacios fundamentales.} La SVD revela la
geometría de $\bm{A}$:
\begin{center}
\begin{tabular}{ll}
\toprule
\textbf{Subespacio} & \textbf{Base ortonormal} \\
\midrule
Espacio columna $\mathrm{col}(\bm{A})$ & $\{\bm{u}_1,\ldots,\bm{u}_r\}$ \\
Núcleo izquierdo $\ker(\bm{A}^\top)$ & $\{\bm{u}_{r+1},\ldots,\bm{u}_m\}$ \\
Espacio fila $\mathrm{fila}(\bm{A})$ & $\{\bm{v}_1,\ldots,\bm{v}_r\}$ \\
Núcleo $\ker(\bm{A})$ & $\{\bm{v}_{r+1},\ldots,\bm{v}_n\}$ \\
\bottomrule
\end{tabular}
\end{center}

\subsection{La pseudoinversa}
\label{subsec:ap_pseudoinversa}

La pseudoinversa (Moore-Penrose) de $\bm{A}$ es:
\begin{equation}
  \bm{A}^\dagger = \bm{V}\bm{\Sigma}^\dagger\bm{U}^\top,
  \qquad
  \bm{\Sigma}^\dagger = \mathrm{diag}\!\left(\frac{1}{\sigma_1},\ldots,
  \frac{1}{\sigma_r},0,\ldots,0\right).
  \label{eq:ap_pseudoinversa}
\end{equation}
Es la inversa generalizada que existe aunque $\bm{A}$ sea rectangular o
deficiente de rango. Sus propiedades fundamentales:
\begin{itemize}
  \item Si $\bm{A}$ tiene rango completo de columnas ($r = n$, sistema
        sobredeterminado): $\bm{A}^\dagger = (\bm{A}^\top\bm{A})^{-1}\bm{A}^\top$
        (inversa izquierda). El estimador $\hat{\bm{\theta}} = \bm{A}^\dagger\bm{b}$
        es la solución de mínimos cuadrados.
  \item Si $\bm{A}$ tiene rango completo de filas ($r = m$, sistema subdeterminado):
        $\bm{A}^\dagger = \bm{A}^\top(\bm{A}\bm{A}^\top)^{-1}$ (inversa derecha).
        $\hat{\bm{\theta}} = \bm{A}^\dagger\bm{b}$ es la solución de mínima norma.
  \item $\bm{A}\bm{A}^\dagger$ es la proyección ortogonal sobre $\mathrm{col}(\bm{A})$.
  \item $\bm{A}^\dagger\bm{A}$ es la proyección ortogonal sobre $\mathrm{fila}(\bm{A})$.
\end{itemize}

\begin{observacion}{La pseudoinversa y el doble descenso}
En el Capítulo~\ref{ch:cap8}, el régimen sobreparamétrico ($p > N$) corresponde
al caso de rango completo de filas: el sistema $\bm{X}\bm{\theta} = \bm{y}$
tiene infinitas soluciones. El descenso de gradiente desde $\bm{\theta}^{(0)} =
\bm{0}$ converge a $\hat{\bm{\theta}} = \bm{X}^\dagger\bm{y}$, la solución de
mínima norma, por la Proposición~\ref{prop:min_norma}. La fórmula del error de
test~\eqref{eq:dd_formula} depende críticamente de que esta sea la solución
seleccionada y no otra de las infinitas posibles.
\end{observacion}

\subsection{Número de condición y sensibilidad numérica}
\label{subsec:ap_condicion}

El \textbf{número de condición} de $\bm{A}$ es:
\begin{equation}
  \kappa(\bm{A}) = \frac{\sigma_{\max}}{\sigma_{\min}} = \sigma_1/\sigma_r.
  \label{eq:ap_condicion}
\end{equation}

Mide la amplificación máxima de errores en la entrada al resolver el sistema
$\bm{A}\bm{\theta} = \bm{b}$: si $\bm{b}$ tiene un error relativo $\varepsilon$,
la solución puede tener un error relativo de hasta $\kappa(\bm{A})\varepsilon$.
Para $\kappa \gg 1$ (sistema mal condicionado), pequeñas perturbaciones del
lado derecho producen grandes cambios en la solución: el problema es
numéricamente inestable.

\medskip\noindent\textbf{Relación con la regularización.} La regularización de
Tikhonov del Capítulo~\ref{ch:cap2} con parámetro $\lambda$ efectivamente
reemplaza $\bm{A}$ por $\bm{A}_\lambda$ con valores singulares modificados
$\sigma_i^{(\lambda)} = \sigma_i^2/(\sigma_i^2 + \lambda)$ (filtro espectral).
El número de condición efectivo es:
\[
  \kappa(\bm{A}_\lambda) \approx \frac{\sigma_1^2}{\sigma_r^2 + \lambda},
\]
que decrece con $\lambda$: la regularización mejora el condicionamiento a costa
de sesgar la solución.

%───────────────────────────────────────────────────────────────────────────────
\section{Normas matriciales}
\label{sec:ap_normas}
%───────────────────────────────────────────────────────────────────────────────

\textit{Usado en: Capítulo~\ref{ch:cap3} (paisaje de pérdida, curvaturas),
Capítulo~\ref{ch:cap8} (cota de Bartlett, gradient clipping, normas durante el
entrenamiento), Capítulo~\ref{ch:cap9} (Lema de Schur, capas equivariantes).}

\medskip
Para una matriz $\bm{W}\in\mathbb{R}^{m\times n}$ con valores singulares
$\sigma_1\geq\cdots\geq\sigma_r > 0$:

\begin{center}
\begin{tabular}{llll}
\toprule
\textbf{Norma} & \textbf{Definición} & \textbf{Expresión en SVD} & \textbf{Rol en el libro} \\
\midrule
Espectral $\|\cdot\|_2$ & $\max_{\|\bm{x}\|=1}\|\bm{W}\bm{x}\|$ & $\sigma_1$ & Cota de Bartlett \\
Frobenius $\|\cdot\|_F$ & $\sqrt{\sum_{ij}W_{ij}^2}$ & $\sqrt{\sum_i\sigma_i^2}$ & Regularización $\ell_2$ \\
Nuclear $\|\cdot\|_*$ & $\sum_i\sigma_i$ & $\sum_i\sigma_i$ & Complejidad efectiva \\
\bottomrule
\end{tabular}
\end{center}

Las tres normas están relacionadas: para $\bm{W}\in\mathbb{R}^{m\times n}$
con rango $r$,
\begin{equation}
  \|\bm{W}\|_2 \leq \|\bm{W}\|_F \leq \sqrt{r}\|\bm{W}\|_2,
  \qquad
  \|\bm{W}\|_2 \leq \|\bm{W}\|_* \leq \sqrt{r}\|\bm{W}\|_F.
  \label{eq:ap_normas_rels}
\end{equation}

\medskip\noindent\textbf{La norma espectral como amplificación máxima.}
$\|\bm{W}\|_2 = \sigma_1$ es la máxima amplificación que $\bm{W}$ puede
impartir a un vector: $\|\bm{W}\bm{x}\|_2\leq\|\bm{W}\|_2\|\bm{x}\|_2$.
En el contexto del retroceso del gradiente (Capítulo~\ref{ch:cap6}), el
gradiente en la capa $l$ es multiplicado por $\bm{W}^{(l)\top}$, cuya norma
espectral es $\sigma_1(\bm{W}^{(l)})$. Si $\|\bm{W}^{(l)}\|_2 < 1$ para
todas las capas, el gradiente se desvanece exponencialmente; si $\|\bm{W}^{(l)}\|_2 > 1$,
explota. Las conexiones residuales del Capítulo~\ref{ch:cap9} resuelven
esto añadiendo un camino de gradiente de norma 1.

\medskip\noindent\textbf{La norma de Frobenius y la regularización $\ell_2$.}
El weight decay $\lambda\sum_l\|\bm{W}^{(l)}\|_F^2$ penaliza la norma de
Frobenius de cada capa. Por la relación~\eqref{eq:ap_normas_rels}, controlar
$\|\bm{W}\|_F$ también controla $\|\bm{W}\|_2$ (aunque de forma menos ajustada).
Las cotas de generalización del Capítulo~\ref{ch:cap8} (Bartlett et al., 2017)
usan la combinación $\prod_l\|\bm{W}^{(l)}\|_2$ y $\sum_l\|\bm{W}^{(l)}\|_F^2/
\|\bm{W}^{(l)}\|_2^2$: dependen de ambas normas simultáneamente.

\medskip\noindent\textbf{Submultiplicatividad.} Para la norma espectral y la
de Frobenius:
\begin{equation}
  \|\bm{A}\bm{B}\|_2 \leq \|\bm{A}\|_2\|\bm{B}\|_2,
  \qquad
  \|\bm{A}\bm{B}\|_F \leq \|\bm{A}\|_2\|\bm{B}\|_F.
  \label{eq:ap_submulti}
\end{equation}
La primera es la razón por la que la cota de Bartlett tiene la forma de un
producto de normas espectrales: la composición de capas lineal no puede ampliar
más que el producto de las amplificaciones de cada capa.

%───────────────────────────────────────────────────────────────────────────────
\section{El hessiano y la curvatura del paisaje de pérdida}
\label{sec:ap_hessiano}
%───────────────────────────────────────────────────────────────────────────────

\textit{Usado en: Capítulo~\ref{ch:cap3} (curvatura, mínimos planos vs.\
estrechos), Capítulo~\ref{ch:cap5} (estabilidad de Euler, temperatura efectiva
del SGD), Capítulo~\ref{ch:cap8} (double descent, generalización).}

\medskip
Sea $f:\mathbb{R}^p\to\mathbb{R}$ de clase $C^2$. La \textbf{matriz hessiana}
en $\bm{\theta}$ es:
\begin{equation}
  \bm{H}(\bm{\theta}) = \nabla^2 f(\bm{\theta}),
  \qquad
  H_{ij} = \frac{\partial^2 f}{\partial\theta_i\partial\theta_j}.
  \label{eq:ap_hessiano}
\end{equation}
$\bm{H}$ es simétrica (si $f\in C^2$, por el teorema de Schwarz) y semidefinida
positiva en un mínimo local.

\medskip\noindent\textbf{Aproximación de Taylor de segundo orden.} En una
vecindad de $\bm{\theta}_0$:
\begin{equation}
  f(\bm{\theta}_0 + \bm{d}) \approx f(\bm{\theta}_0) + \nabla f(\bm{\theta}_0)^\top\bm{d}
  + \frac{1}{2}\bm{d}^\top\bm{H}(\bm{\theta}_0)\bm{d}.
  \label{eq:ap_taylor2}
\end{equation}
Los autovalores del hessiano $\{\lambda_i\}$ determinan la curvatura en las
direcciones de los autovectores correspondientes. En un mínimo ($\nabla f = \bm{0}$):
\begin{itemize}
  \item $\lambda_i > 0$ para todo $i$: mínimo local estricto, paisaje localmente
        cuadrático.
  \item Algunos $\lambda_i = 0$: mínimo degenerado, dirección plana (la pérdida
        no crece en esa dirección).
  \item $\lambda_i < 0$ para algún $i$: punto de silla, no es un mínimo.
\end{itemize}

\subsection{Número de condición del hessiano y el paso de Euler}
\label{subsec:ap_condicion_hessiano}

El número de condición del hessiano $\kappa(\bm{H}) = \lambda_{\max}/\lambda_{\min}$
determina la dificultad de optimización. La estabilidad del método de Euler
(descenso de gradiente) requiere tasa de aprendizaje:
\begin{equation}
  \eta < \frac{2}{\lambda_{\max}(\bm{H})}.
  \label{eq:ap_euler_estabilidad}
\end{equation}
Para un hessiano mal condicionado ($\kappa \gg 1$), la tasa debe ser pequeña
(dictada por $\lambda_{\max}$) pero el progreso en las direcciones de baja
curvatura es lento (proporcional a $\eta\lambda_{\min}$). El número de
iteraciones para converger escala como $O(\kappa)$: el condicionamiento es
la barrera fundamental de la optimización de primer orden.

La normalización por lotes (Capítulo~\ref{ch:cap8}) reduce efectivamente $\kappa$
al normalizar las activaciones: hace que el hessiano de la pérdida sea más
isótropo, permitiendo tasas de aprendizaje más grandes.

\subsection{Mínimos planos vs.\ estrechos y la temperatura del SGD}
\label{subsec:ap_minimos_planos}

Dos mínimos con $f(\bm{\theta}_1^*) = f(\bm{\theta}_2^*)$ pero $\kappa(\bm{H}
(\bm{\theta}_1^*)) \ll \kappa(\bm{H}(\bm{\theta}_2^*))$ se comportan de manera
muy distinta en prueba. El volumen de la cuenca de atracción alrededor de
$\bm{\theta}^*$ es proporcional a $1/\sqrt{\det\bm{H}(\bm{\theta}^*)}$: los
mínimos planos tienen cuencas más amplias.

La distribución de Gibbs del SGD (Capítulo~\ref{ch:cap5}) con temperatura
$T_{\mathrm{eff}} = \eta c/2$ asigna masa relativa a las cuencas proporcional
a su volumen efectivo. En la aproximación gaussiana local de la distribución
de Gibbs alrededor de $\bm{\theta}^*$:
\[
  p^*(\bm{\theta})\approx Z^{-1}\exp\!\left(-\frac{\mathcal{L}(\bm{\theta}^*)}{T_{\mathrm{eff}}}
  - \frac{1}{2T_{\mathrm{eff}}}\bm{d}^\top\bm{H}(\bm{\theta}^*)\bm{d}\right),
\]
cuya constante de normalización local es $(2\pi T_{\mathrm{eff}})^{p/2}/
\sqrt{\det\bm{H}(\bm{\theta}^*)}$. Un mínimo plano (pequeño $\det\bm{H}$)
tiene mayor masa que uno estrecho con la misma profundidad.

%───────────────────────────────────────────────────────────────────────────────
\section{El espacio $L^2$ como espacio de Hilbert}
\label{sec:ap_l2}
%───────────────────────────────────────────────────────────────────────────────

\textit{Usado en: Capítulo~\ref{ch:cap10} (construcción de la integral de Itô
como límite $L^2$, isometría de Itô), Capítulo~\ref{ch:cap11} (adjunto $L^2$
del generador, ecuación de Fokker-Planck como adjunto).}

\medskip
El espacio $L^2(\Omega, \mathcal{F}, \mathbb{P})$ consiste en todas las
variables aleatorias $X$ con $\mathbb{E}[X^2] < \infty$, equipado con el
producto interno:
\begin{equation}
  \langle X, Y\rangle_{L^2} = \mathbb{E}[XY],
  \qquad
  \|X\|_{L^2} = \sqrt{\mathbb{E}[X^2]}.
  \label{eq:ap_l2}
\end{equation}
Con este producto interno, $L^2$ es un \textbf{espacio de Hilbert}: completo
respecto a la norma $\|\cdot\|_{L^2}$, lo que garantiza la existencia de
límites de sucesiones de Cauchy.

\medskip\noindent\textbf{Por qué $L^2$ y no casi seguro.} La integral de Itô
se construye como límite $L^2$ de sumas de Riemann estocásticas:
\[
  \int_0^T f_t\,dW_t = \lim_{n\to\infty}\sum_{k=0}^{n-1}f_{t_k}(W_{t_{k+1}}-W_{t_k})
  \quad\text{en }L^2(\mathbb{P}).
\]
La convergencia casi segura sería más fuerte pero más difícil de establecer;
la convergencia $L^2$ es suficiente para los propósitos del libro y la prueba
es directa desde la isometría.

\subsection{La isometría de Itô como isometría de Hilbert}
\label{subsec:ap_isometria}

La \textbf{isometría de Itô} (Teorema~\ref{thm:isometria} del
Capítulo~\ref{ch:cap10}) es la afirmación de que la integral de Itô
$I: \mathcal{L}^2([0,T]\times\Omega) \to L^2(\Omega)$ es una isometría:
\begin{equation}
  \left\|\int_0^T f_t\,dW_t\right\|_{L^2(\Omega)}^2
  = \int_0^T\|f_t\|_{L^2(\Omega)}^2\,dt.
  \label{eq:ap_iso_ito}
\end{equation}
En términos del producto interno: el mapa $f\mapsto\int_0^T f_t\,dW_t$ preserva
la norma, análogamente a cómo una transformación ortogonal preserva la norma
en $\mathbb{R}^n$. Esto es lo que hace la integral de Itô una extensión
natural de la integral de Riemann ordinaria al mundo estocástico: conserva
la ``energía'' del proceso.

\subsection{El adjunto $L^2$ de un operador diferencial}
\label{subsec:ap_adjunto_l2}

\textit{Usado en: Capítulo~\ref{ch:cap6} (backpropagation como operador adjunto),
Capítulo~\ref{ch:cap11} (ecuación de Fokker-Planck como adjunto del generador).}

\medskip
Sea $\mathcal{A}$ un operador diferencial que actúa sobre funciones $f\in C_c^\infty
(\mathbb{R}^d)$. Su \textbf{adjunto $L^2$} $\mathcal{A}^*$ se define por:
\begin{equation}
  \int (\mathcal{A}f)\,g\,d\bm{x} = \int f\,(\mathcal{A}^*g)\,d\bm{x}
  \quad\text{para toda }f,g\in C_c^\infty(\mathbb{R}^d).
  \label{eq:ap_adjunto_def}
\end{equation}
Es la generalización al caso de operadores diferenciales del adjunto de una
matriz: $(\bm{A}\bm{u},\bm{v}) = (\bm{u},\bm{A}^\top\bm{v})$ en $\mathbb{R}^n$.

\begin{proposicion}{Adjunto del generador de difusión}
\label{prop:ap_adjunto_gen}
Sea $\mathcal{A}f = \bm{\mu}\cdot\nabla f + \frac{1}{2}\mathrm{tr}(\bm{D}\nabla^2 f)$
el generador infinitesimal de la SDE $d\bm{X}_t = \bm{\mu}\,dt + \bm{\Sigma}\,d\bm{W}_t$
con $\bm{D} = \bm{\Sigma}\bm{\Sigma}^\top$. Su adjunto $L^2$ es:
\begin{equation}
  \mathcal{A}^*g
  = -\nabla\cdot(\bm{\mu}\,g) + \frac{1}{2}\sum_{ij}\partial_{x_ix_j}(D_{ij}g),
  \label{eq:ap_adjunto_gen}
\end{equation}
que es exactamente el operador de la ecuación de Fokker-Planck~\eqref{eq:fp}.
\end{proposicion}

\begin{proof}
Por integración por partes con soporte compacto:
$\int\bm{\mu}\cdot\nabla f\cdot g\,d\bm{x} = -\int f\,\nabla\cdot(\bm{\mu}g)\,d\bm{x}$.
Para el término difusivo: $\int D_{ij}\partial_{x_ix_j}f\cdot g\,d\bm{x}
= \int f\,\partial_{x_ix_j}(D_{ij}g)\,d\bm{x}$ (integrando por partes dos veces).
Sumando: $\int(\mathcal{A}f)g\,d\bm{x} = \int f(\mathcal{A}^*g)\,d\bm{x}$
con $\mathcal{A}^*$ dado por~\eqref{eq:ap_adjunto_gen}. $\square$
\end{proof}

La ecuación de Fokker-Planck $\partial_t p_t = \mathcal{A}^*p_t$ describe cómo
evoluciona la densidad de probabilidad; la ecuación de Kolmogorov hacia atrás
$\partial_t u + \mathcal{A}u = 0$ describe cómo evoluciona la esperanza
condicional. La dualidad entre ambas es la dualidad operador/adjunto en $L^2$.

\medskip\noindent\textbf{La misma dualidad en el Capítulo~\ref{ch:cap6}.}
En backpropagation, el forward pass aplica la función $f:\bm{a}\mapsto\bm{W}\bm{a}$
y el backward pass aplica la traspuesta $\bm{W}^\top$, que es exactamente el
adjunto respecto al producto interno estándar de $\mathbb{R}^n$. La traspuesta
matricial y el adjunto $L^2$ del generador de difusión son instancias del mismo
concepto algebraico: el adjunto respecto a un producto interno.

%───────────────────────────────────────────────────────────────────────────────
\section{Proyecciones ortogonales y la descomposición espectral}
\label{sec:ap_proyecciones}
%───────────────────────────────────────────────────────────────────────────────

\textit{Usado en: Capítulo~\ref{ch:cap2} (proyección sobre el espacio columna,
solución de mínima norma), Capítulo~\ref{ch:cap8} (régimen sobreparamétrico,
proyección de $\bm{\theta}^*$ sobre $\ker(\bm{X})$).}

\medskip
Una \textbf{proyección ortogonal} sobre un subespacio $V\subseteq\mathbb{R}^n$
es una matriz $\bm{P}$ que satisface: (1) $\bm{P}^2 = \bm{P}$ (idempotente),
(2) $\bm{P}^\top = \bm{P}$ (simétrica), (3) $\mathrm{col}(\bm{P}) = V$.

De la SVD de $\bm{A} = \bm{U}\bm{\Sigma}\bm{V}^\top$ con rango $r$:
\begin{align}
  \bm{P}_{\mathrm{col}} &= \bm{U}_r\bm{U}_r^\top = \bm{A}\bm{A}^\dagger
  \quad\text{(proyección sobre } \mathrm{col}(\bm{A})\text{)}, \\
  \bm{P}_{\mathrm{fila}} &= \bm{V}_r\bm{V}_r^\top = \bm{A}^\dagger\bm{A}
  \quad\text{(proyección sobre } \mathrm{fila}(\bm{A})\text{)},
\end{align}
donde $\bm{U}_r$ y $\bm{V}_r$ son las primeras $r$ columnas de $\bm{U}$ y
$\bm{V}$.

En el Capítulo~\ref{ch:cap8}, la demostración del Teorema~\ref{thm:double_descent}
para $\gamma > 1$ descompone el error del estimador de mínima norma en:
\[
  \hat{\bm{\theta}} - \bm{\theta}^*
  = \underbrace{(\bm{P}_{\mathrm{fila}}-\bm{I})\bm{\theta}^*}_{\text{sesgo}}
  + \underbrace{\bm{X}^\dagger\bm{\varepsilon}}_{\text{varianza}}.
\]
El sesgo es la proyección de $\bm{\theta}^*$ sobre el núcleo de $\bm{X}$
(complemento ortogonal del espacio fila): la parte de la señal verdadera que
no es recuperable porque vive en las direcciones que $\bm{X}$ no observa.

%───────────────────────────────────────────────────────────────────────────────
\begin{notasbib}
La SVD y la pseudoinversa se presentan con rigor en Golub y Van Loan
(1996)~\cite{golub1996matrix}: la referencia estándar de álgebra lineal numérica.
El teorema de Eckart-Young para la mejor aproximación de bajo rango está en el
Capítulo 2 de esa referencia. Para el contexto del aprendizaje automático, una
presentación accesible está en Strang (2019)~\cite{strang2019linear}.

Las normas matriciales y su rol en las cotas de generalización se discuten en
Bartlett et al.\ (2017)~\cite{bartlett2017spectrally}. La submultiplicatividad
de las normas espectrales y de Frobenius es estándar; véase Horn y Johnson
(2013)~\cite{horn2013matrix}.

La construcción del espacio $L^2$ como espacio de Hilbert, la isometría de Itô
y el adjunto $L^2$ del generador se desarrollan en Karatzas y
Shreve~\cite{karatzas1991brownian}, Capítulos 1 y 5. La derivación de la
ecuación de Fokker-Planck como adjunto del generador via integración por partes
(Proposición~\ref{prop:ap_adjunto_gen}) sigue la presentación de Øksendal
(2003)~\cite{oksendal2003stochastic}, Sección 8.1. La dualidad operador/adjunto
en el contexto de backpropagation se analiza en Griewank y Walther
(2008)~\cite{griewank2008evaluating}.
\end{notasbib}
